% 精读章: Burke, Hsiang & Miguel, Global Non-Linear Effect of Temperature
%         on Economic Production (Nature 2015)
% 原文出处: doi:10.1038/nature15725（所据版本：期刊已刊版 Nature 527: 235--239；
%          本仓 papers/01_methods/2015-Nature-Burke-Hsiang-Miguel-*.pdf 为团队正典，
%          含正文 5 页与 Extended Data Fig. 1--6、Extended Data Table 1--3；
%          本章页码引注均指该版。Supplementary Information（下称 SI）2026-11-19 自
%          作者网站下载，存 papers/01_methods/2015-Nature-...-Supplementary-Information.pdf，
%          该副本数学字体损坏，公式字母不可辨；引 SI 处写节号）
% 体例注记: 2026-10-21 起，正文中的 terms/expressions 一律用英文原词，
%          中文仅保留句法成分；定义处随文给中文对照并登记英文 \index。
% 建档/更新日期由章首 \chdate 渲染（grep: chdate）
% 主笔: 待认领；起草: Claude（2026-11-04，bjr 审订前稿）
% 核对: 2026-11-19 对照本地 PDF 全章重核——原句引用逐条改回原文字面；
%       改正 0.29（原误作 0.28）、75% 所属情景、预测条数、Figure 1 记号、
%       robustness check 清单、损害函数两点差异等处；数据来源据正文参考文献补齐。
% partinfo 定位: nonlinear panel approach 的代表性应用——fixed effects 加二次
%               温度项使横截面 weather 变异进入 identification，见
%               \chref{ch:info-climate-change-adaptation} 的方法线与速查表。

\chapter{Global Non-Linear Effect of Temperature on Economic Production（Burke, Hsiang \& Miguel，Nature 2015）}
\label{ch:burke-hsiang-miguel-2015}
% 章题过长，页眉用短题（preamble 的 \chaptermark 经 \markboth 排页眉）
\chaptermark{Global Non-Linear Temperature Effects（Burke et al.，2015）}
\chdate{2026-11-04}{2026-11-19}

\section{这篇论文在说什么}

一句话概括：论文用 166 国 1960--2010 的年度面板，以``comparing a country to itself in years when it is exposed to warmer- versus cooler-than-average temperatures''\citep[第~236 页]{burke2015global}的组内比较配合二次温度项，发现``country-level economic production is smooth, non-linear, and concave in temperature (Fig.~2a), with a maximum at 13~$^\circ$C''\citep[第~236 页]{burke2015global}，并报告样本期内的变暖没有改变这一响应——``substantial observed warming over the period apparently did not induce notable adaptation''\citep[第~237 页]{burke2015global}，这一判断由此进入领域。

它在 partinfo 里的坐标：nonlinear panel approach 的代表应用站（\chref{ch:info-climate-change-adaptation} 方法线 panel 线第二站）。它被 \citet{merel2021climate} 列入承认 nonlinear panel proposition 的研究名单（第~1208 页），其面板数据也被 Mérel 与 Gammans 用作经验应用素材（第~1224 页），见 \chref{ch:merel-gammans-2021}。

骨架表（Nature Letter 无编号小节；页码为期刊页码；Extended Data 排在正文之后，无期刊页码）：

\begin{table}[htbp]
\centering\small
\begin{tabular}{@{}p{3.4cm}p{1.4cm}p{8.0cm}@{}}
\toprule
段落功能 & 页码 & 讲什么 \\
\midrule
摘要 & 235 & 非线性、13~$^\circ$C 最优、全球通用；projection：不减缓的变暖使 2100 年全球平均收入降约 23\% \\
问题与理论框架 & 235--236 & micro 高度非线性响应如何加总成 macro 平滑凹响应；式 (1)；两条预测 \\
数据与识别 & 236 & 166 国 1960--2010；一国自身冷暖年份的比较；增长的分解 (1)--(4) \\
主要结果 & 236--237 & 13~$^\circ$C 最优、robustness check、滞后效应、rich/poor、农业与非农业 \\
adaptation 的证据 & 237 & 1960--1989 对 1990--2010 两段子样本；未见明显 adaptation \\
Projection & 237--239 & RCP8.5 与 SSP 情景、收入分布、与 IAM 损害函数对比 \\
Extended Data & 无页码 & Fig.~1--6 与 Table~1--3：各项 robustness check 的图表 \\
\bottomrule
\end{tabular}
\end{table}

\section{预备知识}

读者需备两件。其一，增长率回归的用途：论文估计的是``how economic production changes relative to the previous year—that is, annual economic growth—to purge the data of secular factors in each economy that evolve gradually''\citep[第~236 页]{burke2015global}，即以相对上年的变化为对象，把缓慢演变的国别因素差分掉，再用 fixed effects 与趋势项吸收其余。其二，加总记号（第~236 页）：$f_i(T)$ 是行业 $i$ 中单个生产单元（如一个厂商）的生产率对瞬时（如逐日）温度 $T$ 的响应；对给定的国家、时期与行业，$g_i(T-\bar{T})$ 是全部 unit-hours 在各温度上的分布，以平均温度 $\bar{T}$ 为中心；$m_{i1}$ 与 $m_{i2}$ 分别是低于与高于临界温度阈值的 unit-hours 份额。

\section{数据与测度}

产出数据覆盖 166 国、1960--2010 年（第~236 页），取自 World Bank 的 World Development Indicators 2012（正文参考文献~17）；被解释变量是人均 GDP 对数的年度变化，即增长率（Figure~2a 图注，第~236 页；$N=6{,}584$）。climate 变量是年平均温度与降水，取自 Matsuura 与 Willmott 的 Terrestrial Air Temperature and Precipitation 月度与年度序列 v.~3.01（正文参考文献~18）；Extended Data Table~1 表注写明温度以 $^\circ$C、降水以米计。聚合方式写在原文 SI 的 B.1 节：原始数据是 0.5 度格点的月平均温度与月降水总量，按 2000 年的人口密度加权聚合到国家—年份（权重数据取自 Gridded Population of the World）；robustness check 所用的另一套收入数据是 Penn World Tables 8.0；全样本共 6{,}584 个国家—年份观测。

识别所需的全部控制写在一组分解里：``We deconvolve economic growth to account for: (1) all constant differences between countries, for example, culture or history; (2) all common contemporaneous shocks, for example, global price changes or technological innovations; (3) country-specific quadratic trends in growth rates, which may arise, for example, from changing political institutions or economic policies; and (4) the possibly non-linear effects of annual average temperature and rainfall''\citep[第~236 页]{burke2015global}。作者的选型理由是：``This approach is more reliable than only adjusting for observed variables because it accounts for unobserved time-invariant and time-trending covariates''\citep[第~236 页]{burke2015global}。rich 与 poor 的划分用于分组响应：按中位数分为两组（Figure~2b 图注，第~236 页）。SI 的 C.3 节写明分组依据：一国 1980 年经购买力平价（PPP）调整的人均收入低于全球中位数即归入 poor，取 1980 年是因为多数样本国家从这一年起才有 PPP 数据。

\section{模型设定与识别策略：micro 响应如何加总出 macro 响应}

本节是论文的理论贡献。micro 层面的证据早已表明：劳动供给、劳动生产率、作物单产等对温度``exhibit highly non-linear responses''\citep[第~235 页]{burke2015global}，在 20--30\,$^\circ$C 之间存在阈值。而既有的宏观研究报告的是另一幅图景：富国的总体生产率对温度没有响应，穷国只有线性响应（第~235 页）。论文把疑问表述为：``However, it is unclear how these abrupt declines at the micro level are reflected in coarser macro-level data''\citep[第~235--236 页]{burke2015global}。它的答案是加总：若单个 unit-hour 的生产率损失对其他单元影响有限，一国的总产出就是各行业产出之和，每个行业的产出是对该国该时期全部 unit-hours 的积分（式 (1)，第~236 页）
\begin{equation}
Y(\bar{T}) = \sum_i Y_i(\bar{T}) = \sum_i \int_{-\infty}^{\infty} f_i(T)\, g_i(T-\bar{T})\,\mathrm{d}T .
\label{eq:bhm-agg}
\end{equation}
平均温度 $\bar{T}$ 上升时，高于阈值的份额 $m_{i2}$ 逐渐增大，总产出的损失是渐进的、逐步加重的。式~\ref{eq:bhm-agg} 给出两条预测（第~236 页）。其一，$Y(\bar T)$ 是 $\bar{T}$ 的光滑凹函数，其导数等于各 $f_i(T)$ 的导数按各行业在各日温度上的 unit-hours 数加权的平均；micro 响应在阈值处的折点（原文称 kinked）经分布 $g_i$ 加权后被抹平。其二，$Y(\bar T)$ 的峰值所在温度低于 $f_i(T)$ 的阈值，但这是有条件的：``if the slope of $f_i(T)$ above the threshold is steeper than minus the slope below the threshold''\citep[第~236 页]{burke2015global}，作者称 micro 证据支持这一条件。Figure~1（第~235 页）逐步画了这一映射；其 f 图把加总后的斜率标作 $m_{i1}b_{i1}+m_{i2}b_{i2}$，即阈值两侧 micro 斜率按 unit-hours 份额的加权和。

identification 与 \chref{ch:deschenes-greenstone-2007} 同族，但换了尺度：``In practice, we can approximate this experiment by comparing a country to itself in years when it is exposed to warmer- versus cooler-than-average temperatures due to naturally occurring stochastic atmospheric changes''\citep[第~236 页]{burke2015global}。作者明确拒绝截面比较：``We do not compare output across different countries because such comparisons are probably confounded, distinguishing our approach from cross-sectional studies that attribute differences across countries to their temperatures''\citep[第~236 页]{burke2015global}。响应函数取温度的二次式（Extended Data Table~1 的 Temp. 与 Temp.\ sq.\ 两行；基准列的系数为 0.0127 与 $-0.0005$，对应最优温度 13.06\,$^\circ$C）；估计方程是 SI 的 Supplementary Equation (15)（B.2 节）：
\begin{equation}
\Delta Y_{it} = h(T_{it}) + \lambda_1 P_{it} + \lambda_2 P_{it}^2 + \mu_i + \nu_t + \theta_i t + \theta_{i2} t^2 + \eps_{it},
\label{eq:bhm-reg}
\end{equation}
其中 $i$ 为国家、$t$ 为年份，$\Delta Y_{it}$ 是实际人均 GDP 自然对数的一阶差分，$T_{it}$ 与 $P_{it}$ 是年平均温度与年降水，$\mu_i$ 与 $\nu_t$ 是 country 与 year fixed effects，$\theta_i t+\theta_{i2}t^2$ 是各国自己的二次时间趋势；基准设定中 $h(T_{it})$ 取 $T_{it}$ 的二次式。各项的结构与下标依 SI 原式；希腊字母为本讲义所选，因为本仓所收 SI 副本（作者网站版）的数学字体损坏，字母无法辨认。

\begin{ident}[BHM 增长率回归的可搬用设定]
估计方程即式~\ref{eq:bhm-reg}。被解释变量：人均 GDP 对数的年度变化；回归元：年平均温度及其平方、降水及其平方；控制：country fixed effects、year fixed effects、国家二次时间趋势；标准误在国家层面聚类（以上亦见 Extended Data Table~1 表注）。作者对所用 variation 的概括是：考察各国增长率对本国趋势的偏离，是否与本国温度、降水对各自趋势的偏离非线性相关，并扣除各国共同的冲击（第~236 页）。identification assumption 论文没有写成公式；按这一设定，它是「给定上述控制，一国某年的温度偏离与影响增长的其他因素不相关」（本讲义的归纳）。口径要求：温度为全年平均温度，不是 growing season 的温度；产出取增长率而非水平值。
\end{ident}

\begin{intuition}[为什么加总会压低最优温度]
单个生产单元的响应在阈值处由平转陡。一个国家的 unit-hours 沿温度分布摊开，平均温度每升一点，就有一部分 unit-hours 越过阈值，落到损失斜率更陡的一侧。只要阈值以上的斜率比阈值以下的斜率（取负号后）更陡，总产出就会在平均温度尚未达到 micro 阈值时见顶。所以国家层面的最优温度低于 micro 阈值；论文报告的 13~$^\circ$C 与这一预测一致（第~236 页），但式~\ref{eq:bhm-agg} 本身并不给出 13 这个具体数值。
\end{intuition}

\section{结果与解释}

\textbf{主结果。}``We find country-level economic production is smooth, non-linear, and concave in temperature (Fig.~2a), with a maximum at 13~$^\circ$C, well below the threshold values recovered in micro-level analyses and consistent with predictions from equation (1)''\citep[第~236 页]{burke2015global}。响应的形态：``Cold-country productivity increases as annual temperature increases, until the optimum. Productivity declines gradually with further warming, and this decline accelerates at higher temperatures''\citep[第~236 页]{burke2015global}。

\textbf{Robustness check。}作者列出的检验有（第~236--237 页，图表见 Extended Data Fig.~1 与 Extended Data Table~1）：结果具有全球代表性、不由异常值驱动；允许各国的响应随其变富而改变（温度与平均收入交互）；改用高阶多项式或 restricted cubic splines 刻画温度效应；剔除观测少的国家；剔除主要产油国；剔除中国与美国；加入各大洲逐年的共同冲击（continent$\times$year fixed effects）；放松对增长趋势的假定；加入增长率的多期滞后；改用另一套经济数据（Penn World Tables）。温度冲击的滞后效应也考察过，结论是损失之后没有回补：``This `no catch up' behaviour is consistent with the observed response to other climatological disturbances, such as tropical cyclones''\citep[第~237 页]{burke2015global}。

\textbf{rich/poor 与部门分解。}穷国的响应更大——``Poor tropical countries exhibit larger responses mainly because they are hotter on average, not because they are poorer''\citep[第~237 页]{burke2015global}。有迹象显示富国受温度的影响可能略小，但``their response is statistically indistinguishable from poor countries at all temperatures''\citep[第~237 页]{burke2015global}。富国中高温国家少，高温的总效应估计得不精确；即便如此，作者估计有``an 80\% likelihood that the marginal effect of warming is negative at high temperatures in these countries''\citep[第~237 页]{burke2015global}。农业与非农业的总产出对年平均温度都呈非线性，rich 与 poor 两组皆然（第~237 页；Figure~2d、e 在第~236 页）；作者同时说明，低温端在这些子样本里没有显著效应，高温端的负效应在穷国各情形下显著，在富国显著或边际显著（第~237 页）。

\textbf{adaptation 的经验判断。}这是论文进入方法线争论的关键证据：把样本切成 1960--1989 与 1990--2010 两段，``Results using data from 1960–1989 and 1990–2010 are nearly identical (Fig.~2c)''\citep[第~237 页]{burke2015global}；作者的表述是``We do not find that technological advances or the accumulation of wealth and experience since 1960 has fundamentally altered the relationship between productivity and temperature''\citep[第~237 页]{burke2015global}，并进一步说``We find only weak suggestive evidence that richer populations are less vulnerable to warming, and no evidence that experience with high temperatures or technological advances since 1960 have altered the global response to temperature. This suggests that adaptation to climatic change may be more difficult than previously believed''\citep[第~237 页]{burke2015global}。\textbf{注意这条证据的口径}：它是两段子样本之间的系数对比，不是 adaptation 的直接测度；它检验的是「全球响应函数在两段时期之间是否移动」。而按 \chref{ch:merel-gammans-2021} 的分析，这条由二次面板估计出的响应函数本身已经是 short-run 与 long-run response 的加权平均。读这一结论时两层限定要叠加考虑。

\textbf{Projection。}把估计的非线性响应代入 RCP8.5 的变暖情景与 SSP 的基线增长情景（第~237 页）。各国的升温幅度取 CMIP5 全部模式在 RCP8.5 下的集合平均，按人口加权聚合到国家，样本国家到 2100 年的升温在 2.7 到 5.8\,$^\circ$C 之间（SI 的 D.1 节）。基准模型下，``climate change reduces projected global output by 23\% in 2100 (best estimate, SSP5) relative to a world without climate change, although statistical uncertainty allows for positive impacts with probability 0.29''\citep[第~238 页]{burke2015global}。收入分布的后果更重，同样在基准估计下，``average income in the poorest 40\% of countries declines 75\% by 2100 relative to a world without climate change, while the richest 20\% experience slight gains''\citep[第~238 页]{burke2015global}。projection 的前提由作者自己交代：``This approach assumes future economies respond to temperature changes similarly to today's economies—perhaps a reasonable assumption given the observed lack of adaptation during our 50-year sample''\citep[第~237 页]{burke2015global}。与 IAM 损害函数的对比，摘要的概括是：``expected global losses are approximately linear in global mean temperature, with median losses many times larger than leading models indicate''\citep[第~235 页]{burke2015global}。

\section{机制与讨论}

加总机制（式~\ref{eq:bhm-agg}）在结果端的对应是：macro 最优温度 13~$^\circ$C 远低于 micro 研究得到的阈值（20--30\,$^\circ$C），作者称这与式 (1) 的预测一致（第~235--236 页）。不平等机制的表述：``Strong negative correlation between baseline income and baseline temperature indicates that warming may amplify global inequality because hot, poor countries will probably suffer the largest reduction in growth''\citep[第~238 页]{burke2015global}。对 IAM 的含义：作者用自己的结果构造经验损害函数，与 IAM 中常用的损害函数比较，``Commonly used functions are within our estimated uncertainty, but differ in two important respects''\citep[第~238 页]{burke2015global}。两点差异写在第~239 页：其一，论文预测的全球损失对温度大致呈线性（略凹），不是此前理论设定的二次或指数形式；其二，损害函数的斜率在小幅升温时已经很大，2\,$^\circ$C 升温的预期损失是既有估计的 2.5 到 100 倍，更高升温下至少是 2.5 倍。

\section{局限与批评}

论文自己承认的：其一，projection 假定未来经济体对温度的响应与今天相同，其辩护是样本期内未见 adaptation（第~237 页）；作者在结尾处两面都说了，前所未有的创新或防御性投资可能减轻这些影响，社会冲突或贸易受阻则可能加重（第~239 页）。其二，rich 与 poor 的差异估计得不精确：高温富国少，富国高温效应的不确定性大；作者能说的只是``the non-linearity of the rich-country response alone is statistically significant''\citep[第~237 页]{burke2015global}（$P<0.1$）与上文 80\% 的可能性，区分 rich 与 poor 响应的 projection 也因每组数据更少而更不确定（第~237 页）。其三，估计只含温度效应，不含 climate change 的其他损失来源，如热带气旋与海平面上升（第~239 页）。其四，统计不确定性不小，作者处处给出显式的概率（如基准估计下全球影响为正的概率 0.29，第~238 页）。

后续文献指出的：\citet{merel2021climate} 的偏差分解为解读论文的系数补上了缺失的一环——二次面板的温度系数是 long-run 与 short-run response 的加权平均（无 skewness 时为 convex combination），long-run 的权重取决于总体 weather variation 中来自横截面 climatic differences 的份额（第~1209 页）。这套国家面板的横截面差异大，long-run 成分的占比应当较高，但具体份额须按 \chref{ch:merel-gammans-2021} 的权重公式计算后才可断言（本讲义的引申；Mérel 与 Gammans 对全球面板的计算在其在线附录 G.4，本仓未收，待补）。对 projection 的外部批评（收入弹性、加总口径等）在后续文献中存在，待精读后补。

本讲义的评注：论文与 \chref{ch:deschenes-greenstone-2007} 的对照是方法线上最适合放在一起讲的一组——同一族组内 identification，两者的回归都含 weather 变量的二次项（D-G 第~367 页；Mérel 与 Gammans 把两篇同列于二次设定面板的例子，第~1208 页脚注~3），差别在尺度与被解释变量：D-G 用美国县级农业利润，BHM 用国家人均 GDP 增长率。BHM 自己引以为旁证的是另一项美国县级研究：``Our finding that rich countries respond non-linearly to temperature is consistent with recent county-level results in the United States''\citep[第~237 页；所引为 Deryugina 与 Hsiang 2014 的工作稿，正文参考文献~8]{burke2015global}。

\section{在本领域脉络中的位置}

论文回答的旧问题：温度与产出的 macro 关系是否非线性、全球是否共用一条响应（答案：是，13~$^\circ$C 最优）。它并非 nonlinear panel approach 的首创：penalty 结构与「hybrid estimator」最早见于 \citet{mcintosh2006identifying}，农业端的二次 weather 面板至少可追溯到 \citet{deschenes2007economic}（第~367 页；名单见 \citealp[第~1208 页脚注~3]{merel2021climate}）；论文的地位是全球宏观尺度上的旗舰应用。它开的新问题：这条响应里有多少 long-run adaptation 成分——由 \chref{ch:merel-gammans-2021} 形式化为可计算量；「样本期内未见 adaptation」的判断也把 adaptation 的检验推到台前（\chref{ch:deschenes-greenstone-2007} 的自认局限、\chref{ch:hsiang-2016} 对各时间尺度估计的比较，处理的是同一争论的不同侧面）。本章在 \chref{ch:info-climate-change-adaptation} 的方法线与速查表中均有登记。

% 章末参考文献（只列本章引用的条目），勿删
\chapterrefs
